% ============================================================
% EECSI 2026 — IEEE Conference Paper
% RAJDOLL: Multi-Agent LLM Orchestration via MCP for
%           Automated OWASP WSTG 4.2 Web Penetration Testing
%
% Upload to Overleaf → compile with pdfLaTeX
% Template: IEEEtran (conference mode)
% ============================================================
\documentclass[conference]{IEEEtran}

% ---- Packages -----------------------------------------------
\usepackage[T1]{fontenc}
\usepackage[utf8]{inputenc}
\usepackage{cite}
\usepackage{amsmath,amssymb}
\usepackage{graphicx}
\usepackage{booktabs}
\usepackage{array}
\usepackage{url}
\usepackage{hyperref}
\usepackage{microtype}
\usepackage{pifont}
\usepackage{threeparttable}
\newcommand{\cmark}{\ding{51}}
\newcommand{\xmark}{\ding{55}}

% Tighter table columns
\newcolumntype{L}[1]{>{\raggedright\arraybackslash}p{#1}}
\newcolumntype{C}[1]{>{\centering\arraybackslash}p{#1}}
\newcolumntype{R}[1]{>{\raggedleft\arraybackslash}p{#1}}

% ---- Title --------------------------------------------------
\title{RAJDOLL: Multi-Agent LLM Orchestration via Model Context
Protocol for Automated OWASP WSTG~4.2 Web Penetration Testing}

% ---- Author -------------------------------------------------
\author{
  \IEEEauthorblockN{Martua Raja Doli Pangaribuan}
  \IEEEauthorblockA{
    Department of Cyber Security Engineering\\
    Politeknik Siber dan Sandi Negara\\
    Bogor, Indonesia\\
    martua.raja@pssn.ac.id
  }
}

% =============================================================
\begin{document}

\maketitle

% ---- Abstract -----------------------------------------------
\begin{abstract}
Web application penetration testing remains costly and expertise-bound, despite automated scanners and the recent wave of LLM-powered agents. Existing LLM-agent pentest systems rely almost exclusively on cloud LLMs (e.g., GPT-4), raising data-privacy and API-cost concerns, and are rarely evaluated against structured web-security standards. We present RAJDOLL, a local-LLM multi-agent system for web penetration testing that coordinates 14 specialized agents through a Planner-Summarizer Sequential pattern, abstracts security tooling via the Model Context Protocol (MCP), and runs entirely on a consumer-grade Qwen 3-4B model without external API dependencies. Each agent is aligned to a distinct OWASP Web Security Testing Guide~(WSTG) v4.2 category; the Planner-Summarizer pipeline carries a cumulative compressed summary across agents to prevent context loss. We evaluate RAJDOLL on OWASP Juice Shop, achieving Precision~94.87\%, Recall~91.23\%, and F1-Score~93.01\% (TCR~85.19\%), exceeding all acceptance thresholds and completing a full 14-agent scan in under one hour on consumer hardware. To the best of our knowledge, RAJDOLL is the first local-LLM multi-agent pentest system employing MCP for tool abstraction and systematically evaluated against the full OWASP WSTG~4.2 catalog.
\end{abstract}

\begin{IEEEkeywords}
Model Context Protocol, Large Language Model, Multi-Agent System,
Web Penetration Testing, OWASP WSTG 4.2, Local LLM
\end{IEEEkeywords}

% =============================================================
\section{Introduction}

Modern web applications form the entry point for a significant share of real-world compromises, and web penetration testing remains the dominant empirical technique for surfacing pre-deployment risk~\cite{bssn2024}\cite{hackerone2024}. Manual testing is thorough but expensive; automated scanners such as Burp Suite, OWASP ZAP, and Acunetix are faster but weak on context-dependent vulnerabilities. The asymmetry between attacker creativity and scanner coverage persists.

A recent line of work addresses this gap with LLM-agent pentest systems --- PentestGPT~\cite{pentestgpt}, HackSynth~\cite{hacksynth}, and PENTEST-AI~\cite{pentestai} among others. These systems delegate reasoning to a frontier cloud LLM (typically GPT-4) and orchestrate tool invocation through a single agent or ad-hoc multi-agent pipeline. Accuracy improves, but three limitations remain under-examined: the reliance on cloud LLMs for what is fundamentally a privacy-sensitive workload, the absence of a standardized tool-abstraction layer across agents, and the lack of systematic alignment between evaluation metrics and a structured security standard such as OWASP WSTG.

This paper addresses those gaps. We present RAJDOLL, a multi-agent system that runs 14 specialized pentest agents under a Planner-Summarizer Sequential coordinator~\cite{hacksynth}, abstracts security tooling through the Model Context Protocol (MCP)~\cite{mcp}, and uses a local Qwen 3-4B model for all reasoning --- removing external-API dependencies entirely. We evaluate RAJDOLL on OWASP Juice Shop using per-ground-truth-item metrics aligned to OWASP WSTG~v4.2 categories.

Our contributions are fourfold. (i)~\textit{Architectural}: the Planner-Summarizer Sequential pattern for multi-agent LLM coordination with cumulative context, addressing the LLM-contention pathology of naive parallel execution. (ii)~\textit{Systems}: an MCP-based tool-abstraction layer supporting 14 specialized agents and 100+ security tools through containerized JSON-RPC~2.0 endpoints. (iii)~\textit{Deployment}: to the best of our knowledge, the first local-LLM web-pentest agent, eliminating cloud-API dependence and runnable on consumer hardware. (iv)~\textit{Empirical}: vulnerability-coverage metrics (Precision, Recall, F1, TCR) aligned to all 96 WSTG~v4.2 test cases, demonstrating that a local-LLM agent achieves $\geq$90\% Precision and $\geq$90\% Recall on a complex modern SPA benchmark and establishing a replicable per-WSTG-category evaluation methodology.

% =============================================================
\section{Related Work}

\subsection{LLM-based Penetration Testing Agents}

The use of large language models for offensive security has moved from exploratory demos to benchmarked systems. PentestGPT~\cite{pentestgpt} introduced an LLM-driven reasoning loop that issues shell commands against a target and interprets responses to progress through a reconnaissance-exploit-postexploit cycle; accuracy is strong but the system depends on GPT-4 via API. Task completion improves by 228.6\% over GPT-3.5, yet the system scales poorly --- raw tool output accumulates until it saturates available context, with no compression mechanism to preserve signal across phase boundaries.

HackSynth~\cite{hacksynth} directly addresses this accumulation problem with a Planner-Summarizer dual module. The Planner proposes candidate actions; the Summarizer compresses tool output into reusable state. Evaluated on PicoCTF and OverTheWire CTF challenges, the design manages context effectively. What it does not provide is alignment to any structured web-security methodology --- nor a standardized interface for the underlying tools, which limits composability across specialized agents.

PENTEST-AI~\cite{pentestai} brings four specialized agents --- scanning, exploitation, validation, and reporting --- backed by MITRE ATT\&CK knowledge. Designed for general network targets, it lacks both OWASP WSTG alignment and a standardized inter-agent protocol. All three systems assume a cloud LLM (typically GPT-4), creating data-privacy and cost barriers for privacy-sensitive engagements. RAJDOLL inherits HackSynth's Planner-Summarizer design, extends it to 14~agents aligned with WSTG v4.2 categories, and addresses the cloud-LLM barrier by running a local Qwen 3-4B model.

\subsection{Evaluation Methodologies for Security Agents}

How to measure LLM security agent performance is an open question.
Two evaluation paradigms exist. The \textit{CTF-completion} paradigm
measures whether an agent flags or exploits a challenge: NYU CTF
Bench~\cite{nyuctfbench} curates 200 CSAW CTF challenges; Cybench~\cite{cybench}
draws 40 HackTheBox tasks; InterCode~\cite{intercodectf} provides
execution-feedback loops for interactive measurement. These benchmarks
yield binary success rates but lack alignment to any structured security
methodology, making it difficult to diagnose which \textit{test category}
an agent fails. The \textit{vulnerability-coverage} paradigm, used in
this paper, measures how many items in a structured ground truth catalog
(aligned to OWASP WSTG categories) the agent detects, producing
Precision, Recall, and F1 scores that are both auditable and
category-traceable. RAJDOLL adopts the vulnerability-coverage paradigm,
computing per-ground-truth-item metrics against three intentionally
vulnerable targets with ground truth catalogs aligned to WSTG~v4.2.

\subsection{Model Context Protocol and Tool Abstraction}

The Model Context Protocol~\cite{mcp} is a vendor-neutral JSON-RPC specification for exposing tools, resources, and prompts to LLM clients. Originally introduced to standardize tool integration in developer-assistant settings, MCP defines a uniform invocation surface that decouples clients from specific tool implementations. To the best of our knowledge, MCP has not been applied as a tool-abstraction layer in the published pentest-agent literature as of early 2026. RAJDOLL leverages MCP to give 14 specialized agents a shared, uniform interface over 100+ containerized security utilities, enabling each agent to remain focused on its WSTG category without coupling to particular tool implementations.

% =============================================================
\section{System Design}

\subsection{Architecture Overview}

RAJDOLL is a six-layer system, illustrated in
Fig.~\ref{fig:arch}:
\begin{itemize}
  \item \textbf{Presentation Layer}: HTML/CSS/JavaScript
    frontend; FastAPI REST API; WebSocket real-time streaming.
  \item \textbf{LLM API Layer}: Qwen~3-4B deployed locally via
    LM Studio; all inference calls use \texttt{json\_schema}
    response format for deterministic structured output.
  \item \textbf{Orchestration Layer}: LLM Planner (strategy),
    Task Manager (sequential scheduling), Context Coordinator
    (cumulative summary), Result Aggregation (cross-agent
    correlation).
  \item \textbf{MCP Layer}: JSON-RPC~2.0 over HTTP; 14~MCP
    servers acting as tool registries; a PostgreSQL-backed Shared
    Context database persists cross-agent state.
  \item \textbf{Agent Layer}: 14~specialized agents, each
    executing WSTG test cases within its assigned category.
  \item \textbf{Infrastructure Layer}: PostgreSQL, Redis, Celery,
    SQLAlchemy ORM.
\end{itemize}

\begin{figure}[!t]
  \centering
  \includegraphics[width=\columnwidth]{architecture.png}
  \caption{RAJDOLL six-layer system architecture. The
    Orchestration Layer coordinates 14~specialist agents through
    the MCP Layer (JSON-RPC~2.0), which exposes 163~security
    tools as a uniform interface. The Shared Context (PostgreSQL)
    carries the cumulative summary between agents.}
  \label{fig:arch}
\end{figure}

\subsection{Agent-to-WSTG Mapping}

Table~\ref{tab:agents} shows the 14~agents and their WSTG
assignments. WSTG-INPV is deliberately split across two agents:
InputValidationAgent handles SQLi, XSS, SSTI, NoSQL injection,
HTTP parameter pollution, and ReDoS; FileUploadAgent focuses on
upload vulnerabilities, path traversal, and null-byte bypass.
This split avoids a single agent becoming a bottleneck---WSTG-INPV
covers 8 sub-tests spanning very different tool chains (SQLMap, Dalfox,
path-traversal probes)---and increases depth within each subcategory. Crucially,
ReportGenerationAgent always runs last---even if earlier agents
fail---ensuring the scan produces an auditable output regardless
of intermediate errors.

\begin{table}[!t]
\caption{Agent-to-WSTG Mapping}
\label{tab:agents}
\centering
\small
\begin{tabular}{clcc}
\toprule
\textbf{\#} & \textbf{Agent} & \textbf{WSTG Category} & \textbf{Tools} \\
\midrule
1  & ReconnaissanceAgent     & WSTG-INFO        & 28 \\
2  & AuthenticationAgent     & WSTG-ATHN        & 12 \\
3  & SessionManagementAgent  & WSTG-SESS        &  8 \\
4  & InputValidationAgent    & WSTG-INPV        & 24 \\
5  & AuthorizationAgent      & WSTG-ATHZ        & 11 \\
6  & ConfigDeploymentAgent   & WSTG-CONF        & 16 \\
7  & ClientSideAgent         & WSTG-CLNT        & 17 \\
8  & FileUploadAgent         & WSTG-INPV(upload)&  9 \\
9  & APITestingAgent         & WSTG-APIT        &  8 \\
10 & ErrorHandlingAgent      & WSTG-ERRH        &  6 \\
11 & WeakCryptographyAgent   & WSTG-CRYP        &  7 \\
12 & BusinessLogicAgent      & WSTG-BUSL        & 15 \\
13 & IdentityManagementAgent & WSTG-IDNT        &  8 \\
14 & ReportGenerationAgent   & ---              & --- \\
\midrule
   & \textbf{Total}          &                  & \textbf{163} \\
\bottomrule
\end{tabular}
\end{table}

\subsection{Execution Flow}

A scan proceeds through five sequential phases:

\textbf{Phase~1---Reconnaissance}: ReconnaissanceAgent maps
endpoints, tech stack, and JavaScript routes using the Katana
crawler. Findings are written directly to the Shared Context
database, where subsequent agents read them.

\textbf{Phase~1.5---Auto-Login}: The system authenticates using
credentials supplied in the scan request body. The resulting
token (JWT or session cookie) is stored in
\texttt{SharedContext["scan\_credentials"]} and injected into
all subsequent tool calls---eliminating manual session
configuration between agents.

\textbf{Phase~2---LLM Planning}: Based on Phase~1 reconnaissance
findings, the LLM Planner generates a testing strategy that
prioritizes agents, tools, and target endpoints. This planning
step runs under a 5-minute timeout.

\textbf{Phase~3---Sequential Agent Execution}: Agents~2--13 run
one at a time. Before each agent, the orchestrator injects the
full cumulative summary into the agent's planning context. After
execution, the LLM Summarizer condenses that agent's findings
and appends them to the cumulative summary (120-second timeout
per summarization). The growing summary prevents context loss as
the pipeline advances---this is the mechanism, and it works
because compression is explicit rather than hoping the context
window holds.

\textbf{Phase~4---Cross-Agent Analysis}:
\texttt{analyze\_all\_findings()} correlates findings across all
agents, surfacing multi-category attack chains (e.g., SQLi
compounded by an authentication bypass).

\textbf{Phase~5---Report Generation}: ReportGenerationAgent
produces a structured PDF in OWASP WSTG v4.2~\cite{wstg} format using
Jinja2 templates.

\subsection{Per-Agent Workflow: Plan~$\to$~Execute~$\to$~Summarize}

Each agent follows a uniform three-phase loop
(Fig.~\ref{fig:workflow}):
\begin{enumerate}
  \item \textbf{Plan}: The agent reads the cumulative Shared
    Context and asks the LLM to select relevant tools, arguments,
    and target endpoints for its WSTG category.
  \item \textbf{Execute}: Tools are called via MCP
    (JSON-RPC~2.0). Each call passes through: (a)~a circuit
    breaker that halts execution after 5~consecutive failures;
    (b)~LLM argument merging, where LLM-generated arguments
    override tool defaults; (c)~confidence scoring, which rates
    each finding for reliability before persistence.
  \item \textbf{Summarize}: The LLM distills raw tool output
    into a compact summary appended to the Shared Context for
    subsequent agents.
\end{enumerate}

\begin{figure}[!t]
  \centering
  \includegraphics[width=0.85\columnwidth]{architecture.png}
  \caption{Per-agent Plan~$\to$~Execute~$\to$~Summarize loop.
    The LLM reads cumulative Shared Context (Plan), calls MCP
    tools via JSON-RPC~2.0 (Execute), then compresses output
    back into Shared Context (Summarize) for the next agent.}
  \label{fig:workflow}
\end{figure}

This design extends HackSynth's dual-module concept to a
13-agent sequential pipeline with explicit cross-agent shared
state---where each summarization step is a deliberate lossy
compression, not a best-effort approximation.

\subsection{Local-LLM Considerations}

RAJDOLL uses Qwen~3-4B~\cite{qwen3} served locally via LM Studio with an OpenAI-compatible HTTP endpoint. Three design choices emerge from operating at the 4B-parameter scale. First, we deliberately adopt the non-thinking variant of Qwen~3-4B: the thinking variant emits \texttt{<think>} tokens that break \texttt{json\_schema}-enforced structured outputs, which agents rely on for deterministic tool-argument generation. Second, all agent prompts are schema-guarded --- each LLM call specifies a JSON schema the response must conform to, trading some expressive freedom for reliability at this parameter scale. Third, per-call latency of 15--20~s on consumer hardware motivated the sequential rather than parallel execution model: local LLMs contend more aggressively than cloud endpoints under concurrent access, and sequential execution with cumulative summary avoids context thrashing between agents.

The local deployment removes external-API dependence entirely. A complete scan of 14~agents and 100+ tool invocations runs within a self-contained Docker~Compose stack with no LLM network egress. On our reference configuration---LM Studio serving Qwen~3-4B on a consumer workstation (Intel Core~i7, 16~GB RAM) with GPU acceleration, agent orchestration on a second workstation over LAN---a complete OWASP Juice Shop scan finishes in approximately 65~minutes. Deploying both on the same machine adds approximately 10\% scan time overhead from LLM/agent CPU contention.

% =============================================================
\section{Implementation}

RAJDOLL is built on Python~3.11, FastAPI, Celery, PostgreSQL,
Redis, and Docker Compose. All 14~MCP servers run as separate
Docker containers, each discovered by a single generic adapter
(\texttt{mcp\_adapter/server.py}) that exposes Python functions
as MCP tools automatically---no per-tool registration is
required. Integrated tools include SQLMap (SQL injection), Dalfox
(XSS), Nmap (port scanning), Feroxbuster (directory
brute-forcing), WhatWeb (technology identification), Nikto
(server vulnerability scanning), and 120+ custom tools covering
IDOR, CSRF, session management, business logic, and file upload
vulnerabilities.

Development followed Agile SDLC across four iterations:
\begin{itemize}
  \item \textbf{Iteration~1}: Core infrastructure---MCP adapter,
    orchestrator skeleton, PostgreSQL schema,
    ReconnaissanceAgent proof-of-concept (15~initial findings
    on Juice Shop).
  \item \textbf{Iteration~2}: 13~specialist agents---per-agent
    prompts, tool integration, \texttt{json\_schema} enforcement
    for deterministic structured output.
  \item \textbf{Iteration~3}: Full system integration---Planner-Summarizer
    pipeline, auto-login, HITL checkpoints, WebSocket real-time
    updates. First full-pipeline scan (March~2026): Recall~96.5\%
    in 97~minutes.
  \item \textbf{Iteration~4}: Generalization and evaluation
    formalization---corrected a planning defect where agents skipped
    LLM argument generation; fixed cookie-based authentication
    propagation to tool invocations; removed hard-coded
    application-specific endpoint patterns from agent logic. A
    formal 57-entry ground truth database and strict WSTG
    category-prefix matching established the evaluation foundation
    for the validated results in Section~\ref{sec:eval}.
\end{itemize}

The system runs entirely on-premises. No tool output, request
payload, or target hostname leaves the local environment.

% =============================================================
% =============================================================
\section{Evaluation}
\label{sec:eval}

\subsection{Benchmark Setup}

We evaluate RAJDOLL on \textbf{OWASP Juice Shop}~\cite{juiceshop},
a modern single-page application with 57 automatable challenges
spanning 27 WSTG sub-categories. Its challenge list is fully
documented on the official repository, making every ground-truth
entry verifiable. We use 57 ground-truth entries aligned to
OWASP WSTG~v4.2 sub-category codes.

Each target is deployed as a Docker container on the same local
network as RAJDOLL agents, ensuring no external network dependency.
Each target was scanned with ADAPTIVE\_MODE set to \texttt{off}
(all tools execute) and \texttt{HITL\_MODE} set to \texttt{off}
(fully automated). Credentials are injected via the scan API
body (\texttt{POST /api/scans}) and propagated through
\texttt{SharedContext["scan\_credentials"]} to all tool calls.

Hardware: all scans run on a single consumer workstation (Intel
Core~i7, 16~GB RAM) with Qwen~3-4B~\cite{qwen3} served on a
separate machine via LM Studio at 15--20~seconds per inference call.

\subsection{Evaluation Metrics}

\textbf{Unit of analysis.} All OWASP-benchmark metrics are
computed at the \textit{per-ground-truth-item} level: each entry
in the ground truth catalog is either detected or not, and each
non-info-severity finding either matches a catalog entry (TP) or
does not (FP). This eliminates metric mixing between
finding-level Precision and challenge-level Recall.

\textbf{Matching rule.} A finding with WSTG category code $f_c$
matches ground truth entry with code $g_c$ when:
\begin{equation}
f_c = g_c \;\vee\; g_c \text{ starts with } f_c\text{-} \;\vee\; f_c \text{ starts with } g_c\text{-}
\label{eq:match}
\end{equation}
This allows parent-level codes (e.g., \texttt{WSTG-ATHN}) to
match child-level ground truth (e.g., \texttt{WSTG-ATHN-07}),
handling cases where the agent classifies at WSTG category
granularity rather than sub-test granularity.

\textbf{Precision} measures the fraction of non-trivial findings
confirmed by the ground truth catalog:
\begin{equation}
\text{Precision} = \frac{|\{f \in F : \exists g \in GT, \text{matches}(f,g)\}|}{|F_{\text{non-info}}|}
\label{eq:precision}
\end{equation}
On an intentionally-vulnerable target such as Juice Shop,
any finding whose WSTG category matches no catalog entry is
classified FP---because the catalog is exhaustive for each
security level.

\textbf{Recall} measures the fraction of ground truth entries
covered by at least one finding:
\begin{equation}
\text{Recall} = \frac{|\{g \in GT : \exists f \in F, \text{matches}(f,g)\}|}{|GT|}
\label{eq:recall}
\end{equation}

\begin{equation}
\text{F1} = \frac{2 \times \text{Precision} \times \text{Recall}}{\text{Precision} + \text{Recall}}
\label{eq:f1}
\end{equation}

\textbf{Test Case Rate (TCR)} measures the fraction of WSTG
category-level test cases covered by at least one finding,
independently of whether the finding matches a specific ground
truth entry. TCR captures \textit{breadth} of methodology
coverage; Recall captures \textit{depth} (specific
vulnerabilities detected).
\begin{equation}
\text{TCR} = \frac{|\{c \in C_{WSTG} : \exists f \in F, f_c \text{ prefix-matches } c\}|}{|C_{WSTG}|}
\label{eq:tcr}
\end{equation}

\subsection{Experimental Protocol}

Each benchmark scan is triggered via
\texttt{POST /api/scans} with \texttt{whitelist\_domain} set
to the target container name (e.g., \texttt{"juice-shop"}). All
13~testing agents (agents~1--13 in DEFAULT\_PLAN order, excluding
ReportGenerationAgent) run sequentially; ADAPTIVE\_MODE is
\texttt{off} (all 163~tools execute). Credentials are passed in
the request body. Findings are collected via
\texttt{GET /api/scans/\{id\}/findings} after
ReportGenerationAgent completes. Each benchmark target container
is connected to the \texttt{rajdoll\_default} Docker network so
agents resolve hosts by container name (no DNS or network
configuration required). Matching between findings and ground
truth uses Eq.~\eqref{eq:match}; any non-info finding with no
matching ground truth entry is classified a false positive.

\subsection{Results}
\label{sec:results}

Table~\ref{tab:results} reports Precision, Recall, F1-Score, and Test
Case Rate (TCR) for OWASP Juice Shop.

\begin{table}[!t]
\caption{RAJDOLL Evaluation Results on OWASP Juice Shop}
\label{tab:results}
\centering
\small
\begin{tabular}{lC{1.05cm}C{1.05cm}C{1.05cm}C{0.9cm}C{0.6cm}}
\toprule
\textbf{Target} & \textbf{P (\%)} & \textbf{R (\%)} & \textbf{F1 (\%)} & \textbf{TCR (\%)} & \textbf{Pass?} \\
\midrule
Juice Shop  & 94.87 & 91.23 & 93.01 & 85.19  & \cmark \\
\midrule
\multicolumn{6}{l}{\small Threshold: P$\geq$90\%, R$\geq$80\%, F1$\geq$85\%, TCR$\geq$70\%} \\
\bottomrule
\end{tabular}
\end{table}

\textbf{Juice Shop.}
On OWASP Juice Shop (57 ground truth entries across 27 WSTG
sub-categories), RAJDOLL achieves Precision~94.87\%, Recall~91.23\%,
F1-Score~93.01\%, and TCR~85.19\% (23/27 WSTG categories detected),
exceeding all acceptance thresholds. Five ground truth challenges
remain undetected: one SSRF challenge (WSTG-INPV-19) requires
out-of-band callback infrastructure unavailable in the Docker-local
environment; the remaining four require multi-step authenticated-browser
interactions that automated scanners cannot replicate without browser
automation. The precision improvement from 75.9\% to 94.87\% was
achieved by correcting five WSTG category misclassifications
(CONFIG-04$\to$CONF-04, ATHN-06$\to$ATHZ-02, CLNT-15$\to$CLNT-13)
and removing a redundant dual-tag that inflated FP count. All 14 agents
complete with zero failures. A full scan completes in under one hour on
consumer hardware (Intel Core~i7, 16~GB RAM, Qwen~3-4B at
15--20~s per LLM call).

\textbf{Agent completion.}
On Juice Shop, all 14 agents complete with zero failures
(100\% agent completion rate). The circuit-breaker and timeout
design prevents a failing individual tool from blocking the agent,
and a failing agent from blocking the pipeline---a necessary property
for unattended scanning.

\begin{table}[!t]
\caption{Comparison with Prior LLM-Agent Pentest Systems}
\label{tab:comparison}
\centering
\small
\begin{threeparttable}
\begin{tabular}{lcccc}
\toprule
\textbf{System} & \textbf{LLM} & \textbf{Tool IF} & \textbf{WSTG} & \textbf{P/R/F1} \\
\midrule
PentestGPT~\cite{pentestgpt}  & Cloud & \xmark & \xmark & N/A$^\star$ \\
HackSynth~\cite{hacksynth}    & Cloud & \xmark & \xmark & N/A$^\star$ \\
PENTEST-AI~\cite{pentestai}   & Cloud & \xmark & Part. & N/A$^\star$ \\
\textbf{RAJDOLL (ours)}       & \textbf{Local} & \textbf{MCP} & \textbf{Full} & \textbf{95/91/93} \\
\bottomrule
\end{tabular}
\begin{tablenotes}
  \footnotesize
  \item[$^\star$] Reports CTF/task completion; no P/R/F1 on OWASP benchmarks.
  \item Tool~IF~=~tool interface; WSTG~=~OWASP WSTG alignment.
\end{tablenotes}
\end{threeparttable}
\end{table}

% =============================================================
\section{Discussion}

\subsection{Planner-Summarizer Sequential Pattern}

The cumulative summary mechanism resolves PentestGPT's
context-loss problem---but the resolution is not simply
``more context.'' Rather, it is \textit{structured compression}.
Rather than passing raw tool output to subsequent agents (which
saturates the LLM context window within 2--3 agents on typical
engagements), RAJDOLL condenses each agent's findings into a
dense, category-tagged summary before advancing. The effect is
observable on Juice Shop: AuthenticationAgent read
ReconnaissanceAgent's summary, identified the JWT implementation
in the tech stack reported by WhatWeb and Katana, and focused
all 12~of its tools on JWT-specific attacks---producing 10~findings
in under 3~minutes without any target-specific prompt
configuration. No agent received instructions about Juice Shop;
the behavior emerged entirely from the shared context populated
by the preceding agent. This emergent specialization is, we
argue, what distinguishes the Planner-Summarizer Sequential
pattern from simple tool chaining.

\subsection{MCP as a Unifying Interface}

MCP enables RAJDOLL to call SQLMap, Dalfox, and Nmap through
identical JSON-RPC~2.0 invocations---agents require no knowledge
of individual tool CLI syntax. What follows from this design
choice is extensibility: adding a new security tool requires only
one Python function decorated as an MCP tool; the adapter
auto-discovers it. This is not achievable with direct subprocess
calls or REST-per-tool approaches, both of which require
per-agent adaptation code whenever a new tool is integrated.

\subsection{Limitations and Future Directions}

Three limitations constrain the current system. First, five
ground truth challenges remain undetected on Juice Shop---one SSRF
challenge (WSTG-INPV-19) and four Expert-difficulty challenges
requiring manual browser interaction. Adding a dedicated SSRF
probe is estimated to raise Recall from 91.23\% to approximately
93\%. Second, the Planner-Summarizer loop introduces a
dependency on LLM planning latency: at 15--20~seconds per
inference call on the local Qwen~3-4B model, a 14-agent scan
takes approximately one hour; faster hardware or a larger
parameter count would reduce this but is not required to meet the
4-hour scan budget.

Third, RAJDOLL's endpoint discovery is designed for REST/SPA targets
and does not enumerate PHP form parameters or interact with
traditional form-based web applications. Future work will add
form-aware crawling (e.g., Playwright-based interaction) to extend
coverage to PHP form-based targets.

% =============================================================
\section{Conclusion}

RAJDOLL demonstrates that aligning a multi-agent LLM architecture
directly to OWASP WSTG~4.2 categories---connected through the Model
Context Protocol---produces validated detection accuracy
(Precision~94.87\%, Recall~91.23\%, F1-Score~93.01\%, TCR~85.19\%)
on OWASP Juice Shop, completing a full 14-agent scan in under one hour
on consumer hardware. The Planner-Summarizer Sequential pattern resolves
cross-agent context loss without cloud dependency or context-window
expansion: each agent summarizes its findings into a compact,
category-tagged block that the next agent reads---explicit compression,
not approximation.

Future work will target the remaining five undetected Juice Shop
challenges---beginning with a dedicated SSRF probe and
Playwright-based browser interaction (estimated Recall gain:
$+$2\%)---extend the Human-in-the-Loop director mode for controlled
production-environment deployments, and add form-aware crawling to
extend coverage to PHP form-based applications.

% =============================================================
% References
% =============================================================
\begin{thebibliography}{10}

\bibitem{bssn2024}
Badan Siber dan Sandi Negara (BSSN),
``Lanskap Keamanan Siber Indonesia 2024,''
Jakarta, 2024.

\bibitem{hackerone2024}
HackerOne,
``The 2024 Annual Hacker-Powered Security Report,''
HackerOne, San Francisco, 2024.

\bibitem{pentestgpt}
D.~Deng, C.~Zhou, J.~Li, Y.~Zheng, Y.~Liu, and Y.~Liu,
``PentestGPT: Evaluating and Harnessing Large Language Models
for Automated Penetration Testing,''
in \textit{Proc. 33rd USENIX Security Symp.},
Philadelphia, PA, 2024.

\bibitem{hacksynth}
I.~Abramov, A.~Grishina, A.~Kapitonov, E.~Moiseev,
A.~Yakovlev, and A.~Yunusov,
``HackSynth: LLM Agent and Evaluation Framework for Autonomous
Penetration Testing,''
\textit{arXiv preprint arXiv:2412.01778}, 2024.

\bibitem{pentestai}
A.~E.~Ezugwu \textit{et al.},
``PENTEST-AI: A Multi-Agent AI System for Cybersecurity
Testing,''
in \textit{Proc. 2024 IEEE Int. Conf. Advanced Networking and
Telecommunications (ANTEL)}, 2024.

\bibitem{wstg}
OWASP Foundation,
``Web Security Testing Guide v4.2,''
OWASP, 2021.
[Online]. Available: \url{https://owasp.org/www-project-web-security-testing-guide/}

\bibitem{mcp}
Anthropic,
``Model Context Protocol,''
Nov.\ 2024.
[Online]. Available: \url{https://modelcontextprotocol.io/}

\bibitem{juiceshop}
OWASP Foundation,
``OWASP Juice Shop,''
2024.
[Online]. Available: \url{https://owasp.org/www-project-juice-shop/}

\bibitem{nyuctfbench}
M.~Shao, S.~Jancheska, M.~Udeshi, B.~Dolan-Gavitt, H.~Xi, K.~Milner, B.~Chen, M.~Yin, S.~Garg, P.~Krishnamurthy, F.~Khorrami, R.~Karri, and M.~Shafique,
``NYU CTF Bench: A Scalable Open-Source Benchmark Dataset for Evaluating LLMs in Offensive Security,''
\emph{arXiv preprint arXiv:2406.05590}, 2024.

\bibitem{cybench}
A.~K.~Zhang, N.~Perry, R.~Dulepet, J.~Ji, C.~Menders, J.~W.~Lin, E.~Jones, G.~Hussein, S.~Liu, D.~Jasper, P.~Peetathawatchai, A.~Glenn, V.~Sivashankar, D.~Zamoshchin, L.~Glikbarg, D.~Askaryar, M.~Yang, T.~Zhang, R.~Alluri, N.~Tran, R.~Sangpisit, P.~Yiorkadjis, K.~Osele, G.~Raghupathi, D.~Boneh, D.~E.~Ho, and P.~Liang,
``Cybench: A Framework for Evaluating Cybersecurity Capabilities and Risks of Language Models,''
\emph{arXiv preprint arXiv:2408.08926}, 2024.

\bibitem{intercodectf}
J.~Yang, A.~Prabhakar, K.~Narasimhan, and S.~Yao,
``InterCode: Standardizing and Benchmarking Interactive Coding with Execution Feedback,''
in \emph{Advances in Neural Information Processing Systems (NeurIPS)}, 2023.

\bibitem{qwen3}
Qwen Team,
``Qwen3 Technical Report,''
\emph{arXiv preprint arXiv:2505.09388}, 2025.

\end{thebibliography}

\end{document}
